Let $N$ be the nine-point centre, the midpoint of $OH$. Use Euler\textquotesingle s identity $OI^2=R(R-2r)$ and Feuerbach\textquotesingle s identity $I N = \dfrac{R}{2} - r$ for the internal tangency of the incircle and nine-point circle. The median identity in triangle $OIH$ gives $O I^2 + I H^2 = 2 I N^2 + \dfrac{O H^2}{2}$. With $R=12,r=5,IH=4$, this becomes $24 + 16 = 2 \cdot 1^2 + \dfrac{O H^2}{2}$, so $\boxed{OH^2=76}$.
\begin{asy}
import graph;
usepackage("amsmath");
import geometry;
size(7cm);
pair O=(0,0),H=(sqrt(76),0),N=H/2,I=(21/sqrt(19),sqrt(15/19));draw(O--I--H--O);dot(N);label("$O$",O,(-1,0));label("$H$",H,(1,0));label("$N$",N,(0,-1));label("$I$",I,(0,1));label("median identity (schematic)",(4,-1),(0,-1));
\end{asy}

